% =====================================================================
% IV. OBJETIVOS
% =====================================================================
\section{Objetivos}

\subsection{Objetivo General}
\label{sec:objetivo-general}

Dise\~nar e implementar una plataforma de digitalizaci\'on asistida de partituras que integre reconocimiento \'optico de m\'usica, validaci\'on autom\'atica mediante reglas de teor\'ia musical, correcci\'on interactiva y aprendizaje activo, con el fin de reducir el esfuerzo de revisi\'on humana y mejorar la fidelidad de las transcripciones digitales de partituras impresas y manuscritas modernas.

\subsection{Objetivos Espec\'ificos}
\label{sec:objetivos-especificos}

\begin{enumerate}
    \item Analizar las t\'ecnicas de reconocimiento \'optico de m\'usica, validaci\'on basada en reglas musicales y aprendizaje activo aplicables a la digitalizaci\'on de partituras impresas y manuscritas modernas, con el prop\'osito de establecer los fundamentos t\'ecnicos y metodol\'ogicos de la plataforma.

    \item Dise\~nar una arquitectura de software que integre un modelo OMR, un m\'odulo de validaci\'on autom\'atica de errores musicales, una interfaz de correcci\'on asistida y un mecanismo de aprendizaje activo basado en las correcciones realizadas por el usuario.

    \item Implementar un prototipo funcional de la plataforma Cadenza capaz de procesar im\'agenes de partituras, generar transcripciones en formatos simb\'olicos como MusicXML/MIDI, identificar inconsistencias mediante reglas musicales y permitir la correcci\'on interactiva de los resultados.

    \item Desarrollar un mecanismo de aprendizaje activo que utilice las correcciones realizadas por los usuarios para seleccionar muestras relevantes y contribuir al mejoramiento progresivo del reconocimiento de partituras.

    \item Evaluar experimentalmente el desempe\~no del flujo de digitalizaci\'on asistida frente a una l\'inea base OMR no asistida, utilizando conjuntos de datos de referencia y m\'etricas de fidelidad de transcripci\'on y esfuerzo de correcci\'on humano, particularmente la tasa de error simb\'olico (SER), el tiempo de edici\'on y las intervenciones manuales por comp\'as.
\end{enumerate}
